% !TeX root = ../main.tex
% Fonti: Lezione 14, pp. 13--15; Lezione 15, pp. 3, 6--8.
\section{Successioni geometriche}
\label{sec:successioni-geometriche}

\begin{teorema}
Per $q\in\R$ e $n\ge1$, la successione geometrica $(q^n)$ ha il seguente comportamento:
\[
q^n\longrightarrow
\begin{cases}
+\infty,&q>1,\\
1,&q=1,\\
0,&|q|<1,
\end{cases}
\qquad
(q^n)\text{ è irregolare se }q\le-1.
\]
\end{teorema}
\begin{proof}
\passomath{q>1} La disuguaglianza di Bernoulli, già dimostrata nel capitolo~\ref{cap:naturali-combinatoria}, dà
\[
q^n=(1+(q-1))^n\ge1+n(q-1)\longrightarrow+\infty.
\]
Il confronto implica $q^n\to+\infty$.

\passomath{|q|<1} Se $q=0$, tutti i termini sono nulli. Se $0<|q|<1$, allora $1/|q|>1$, dunque $(1/|q|)^n\to+\infty$ e $|q|^n\to0$. Poiché $-|q|^n\le q^n\le|q|^n$, i carabinieri danno $q^n\to0$, anche per $q$ negativo.

\passomath{q=1\text{ oppure }q\le-1} Per $q=1$ la successione è costante. Per $q=-1$ le estratte pari e dispari sono costanti con valori $1$ e $-1$. Per $q<-1$, scrivendo $q^n=(-1)^n|q|^n$, le estratte pari e dispari tendono rispettivamente a $+\infty$ e $-\infty$. Negli ultimi due casi non esiste limite nella retta estesa.
\end{proof}

\section{Il criterio del rapporto}
% Fonte: Lezione 15, p. 3; dimostrazione assegnata per esercizio, qui ricostruita.

\begin{proposizione}[Criterio del rapporto per successioni positive]
\label{prop:criterio-rapporto}
Sia $a_n>0$ definitivamente e supponiamo che esista
\[
\frac{a_{n+1}}{a_n}\longrightarrow\ell\in[0,+\infty].
\]
\begin{enumerate}[label=\arabic*.]
\item Se $0\le\ell<1$, allora $a_n\to0$.
\item Se $1<\ell\le+\infty$, allora $a_n\to+\infty$.
\end{enumerate}
Per $\ell=1$ il criterio non dà una conclusione.
\end{proposizione}
\begin{proof}
\begin{enumerate}[label=\arabic*.]
\item \textbf{Rapporto minore di uno.} Scegliamo $r$ con $\ell<r<1$. Esiste $N$ tale che $0<a_{n+1}/a_n\le r$ per $n\ge N$. Iterando la disuguaglianza,
\[
0<a_n\le a_Nr^{n-N}\qquad(n\ge N).
\]
La maggiorante tende a zero per il teorema sulle geometriche, quindi $a_n\to0$.
\item \textbf{Rapporto maggiore di uno.} Scegliamo $r>1$ con $r<\ell$, o un qualunque $r>1$ se $\ell=+\infty$. Definitivamente $a_{n+1}\ge r a_n$, da cui $a_n\ge a_Nr^{n-N}\to+\infty$.
\end{enumerate}
\end{proof}

\begin{esempio}[Il caso inconcludente]
Le successioni positive $a_n=n$, $b_n=1/n$ e $c_n=1$, per $n\ge1$, hanno tutte rapporto fra termini consecutivi tendente a $1$, ma tendono rispettivamente a $+\infty$, $0$ e $1$. Il limite del rapporto, in questo caso, non determina il limite della successione.
\end{esempio}

\section{Radici di indice variabile}
% Fonte: Lezione 15, pp. 6--7; completati i casi e la prova assegnata per casa.

\begin{proposizione}
\label{prop:radice-costante}
Per ogni $q>0$, $\sqrt[n]{q}\to1$.
\end{proposizione}
\begin{proof}
Se $q>1$, poniamo $u_n=\sqrt[n]{q}-1>0$. Per Bernoulli,
\[
q=(1+u_n)^n\ge1+nu_n,
\qquad 0\le u_n\le\frac{q-1}{n}\longrightarrow0.
\]
I carabinieri danno $u_n\to0$, quindi $\sqrt[n]{q}\to1$. Per $q=1$ la successione è costante; se $0<q<1$, si applica il risultato a $1/q>1$ e poi si prende il reciproco, perché $\sqrt[n]{q}=1/\sqrt[n]{1/q}$.
\end{proof}

\begin{proposizione}[Dal rapporto alla radice]
\label{prop:rapporto-radice}
Se $a_n>0$ definitivamente e $a_{n+1}/a_n\to\ell\in[0,+\infty]$, allora $\sqrt[n]{a_n}\to\ell$.
\end{proposizione}
\begin{proof}
\textbf{Limite positivo finito.} Fissiamo $0<r<\ell<s$. Definitivamente $r\le a_{n+1}/a_n\le s$. Iterando da un indice $N$,
\[
a_Nr^{n-N}\le a_n\le a_Ns^{n-N}.
\]
Prendendo le radici $n$-esime,
\[
r\sqrt[n]{a_Nr^{-N}}\le\sqrt[n]{a_n}\le s\sqrt[n]{a_Ns^{-N}}.
\]
I due fattori sotto radice sono costanti positive, dunque le successioni esterne tendono a $r$ e $s$. Poiché si possono scegliere $r$ e $s$ arbitrariamente vicini a $\ell$, la radice tende a $\ell$: dato $\varepsilon>0$, si scelgono prima $r,s$ dentro $(\ell-\varepsilon,\ell+\varepsilon)$, poi una soglia comune per i due confronti.

\textbf{Limite nullo.} Per ogni $s>0$, definitivamente $a_{n+1}/a_n\le s$, e la stessa iterazione dà $0<\sqrt[n]{a_n}\le s\sqrt[n]{a_Ns^{-N}}$. Per $\varepsilon>0$ si sceglie $s=\varepsilon/2$; il fattore sotto radice tende a $1$, quindi la maggiorante è minore di $\varepsilon$ definitivamente.

\textbf{Limite infinito.} Per ogni $r>0$, definitivamente $a_{n+1}/a_n\ge r$, quindi $\sqrt[n]{a_n}\ge r\sqrt[n]{a_Nr^{-N}}$. Dato $M>0$, scegliamo $r=2M$; il fattore sotto radice è definitivamente maggiore di $1/2$, perciò $\sqrt[n]{a_n}>M$.
\end{proof}

\begin{esempio}
Per $a_n=n$, il rapporto $(n+1)/n\to1$, dunque $\sqrt[n]{n}\to1$. Più in generale, per $a_n=n^k$, con $k>0$ fissato,
\[
\frac{a_{n+1}}{a_n}=\left(1+\frac1n\right)^k\longrightarrow1
\]
per continuità della potenza sul dominio positivo; quindi $\sqrt[n]{n^k}\to1$.
\end{esempio}

\begin{osservazione}[La forma infinito elevato a zero]
Con una base tendente a $+\infty$ e un esponente tendente a zero, il limite non è determinato dai soli limiti dei due elementi. Infatti
\[
n^{1/n}\longrightarrow1,
\qquad
(n^n)^{1/n}=n\longrightarrow+\infty.
\]
Le basi devono essere positive per considerare senza ambiguità le potenze a esponente reale.
\end{osservazione}
